% Part: many-valued-logic
% Chapter: three-valued-logics
% Section: goedel

\documentclass[../../../include/open-logic-section]{subfiles}

\begin{document}

\olfileid{mvl}{thr}{god}

\olsection{منطقات غودل}

قدم كورت غودل متتالية من المنطقات ذات~$n$ من القيم، يحتوي كل منها
جميع~!!{formula}s الصحيحة في المنطق الحدسي، ويحتويه المنطق الكلاسيكي.
وفيما يأتي أول مثال غير تافه:

\begin{defn}\ollabel{defn:goedel}
يعرّف \emph{منطق غودل الثلاثي القيم}~$\LogGod$ بالمصفوفة الآتية:
\begin{enumerate}
  \item لغة القضايا القياسية~$\Lang L_0$ بروابطها
  $\lfalse$ و$\lnot$ و$\land$ و$\lor$ و$\lif$.
  \item مجموعة قيم الصدق~$V = \{\True, \Undef, \False\}$.
  \item $\True$ هي القيمة المميزة الوحيدة؛ أي~$V^+ = \{\True\}$.
  \item أما~$\lfalse$، فلدينا~$\tf{\lfalse} = \False$. وأما دوال
  صدق الروابط الباقية فتعطيها الجداول الآتية:
  \begin{center}
    \begin{tabular}{c|c}
      $\tf{\lnot}[\LogGod]$ & \\
      \hline
      $\True$ & $\False$ \\
      $\Undef$ & $\False$ \\
      $\False$ & $\True$
    \end{tabular}
    \quad
    \begin{tabular}{c|ccc}
      $\tf{\land}[\LogGod]$ & $\True$ & $\Undef$ & $\False$ \\
      \hline
      $\True$ & $\True$ & $\Undef$ & $\False$ \\
      $\Undef$ & $\Undef$ & $\Undef$ & $\False$\\
      $\False$ & $\False$ & $\False$ & $\False$
    \end{tabular}
    \\[2ex]
    \begin{tabular}{c|ccc}
      $\tf{\lor}[\LogGod]$ & $\True$ & $\Undef$ & $\False$ \\
      \hline
      $\True$ & $\True$ & $\True$ & $\True$ \\
      $\Undef$ & $\True$ & $\Undef$ & $\Undef$ \\
      $\False$ & $\True$ & $\Undef$ & $\False$
    \end{tabular}
    \quad
    \begin{tabular}{c|ccc}
      $\tf{\lif}[\LogGod]$ & $\True$ & $\Undef$ & $\False$ \\
      \hline
      $\True$ & $\True$ & $\Undef$ & $\False$ \\
      $\Undef$ & $\True$ & $\True$ & $\False$  \\
      $\False$ & $\True$ & $\True$ & $\True$
    \end{tabular}
  \end{center}
\end{enumerate}
\end{defn}

لاحظ أن جدولي صدق~$\land$ و$\lor$ هما نفسيهما في منطق~\L
ukasiewicz ومنطق كليني القوي، لكن جدولي~$\lnot$ و$\lif$ يختلفان في
كل منها. ففي منطق غودل~$\tf{\lnot}(\Undef) = \False$. وعلى خلاف منطق
\L ukasiewicz ومنطق كليني، يكون
$\tf{\lif}(\Undef, \False) = \False$؛ وعلى خلاف منطق كليني (لكن كما
في منطق~\L ukasiewicz)، يكون
$\tf{\lif}(\Undef, \Undef) = \True$.

وكما يوحي الارتباط بالمنطق الحدسي المشار إليه أعلاه، فإن~$\LogGod[3]$
قريب من المنطق الحدسي. فجميع الحقائق الحدسية تحصيلات منطقية في
$\LogGod[3]$، وكثير من التحصيلات المنطقية الكلاسيكية غير الصحيحة
حدسيًا لا تكون تحصيلات في~$\LogGod[3]$ أيضًا. فمثلًا، ليست الصيغ الآتية
تحصيلات منطقية:
\begin{align*}
  & p \lor \lnot p && (p \lif q) \lif (\lnot p \lor q) \\
  & \lnot\lnot p \lif p && \lnot(\lnot p \land \lnot q) \lif (p \lor q) \\
  & ((p \lif q) \lif p) \lif p && \lnot(p \lif q) \lif (p \land \lnot q)
\end{align*}
غير أن بعض تحصيلات~$\LogGod[3]$ ليست صحيحة حدسيًا، مثل
$\lnot\lnot p \lor \lnot p$ أو~$(p \lif q) \lor (q \lif p)$.

\begin{prob}
  أعط جداول صدق تبين أن الصيغ الآتية تحصيلات منطقية في~$\LogGod[3]$:
  \begin{align*}
    & \lnot\lnot p \lor \lnot p\\
    & (p \lif q) \lor (q \lif p) \\
    & \lnot(p \land q) \lif (\lnot p \lor \lnot q) \\
    & (p \lif q) \lor (q \lif r) \lor (r \lif s)
  \end{align*}
\end{prob}

\begin{prob}
  أعط جداول صدق تبين أن الصيغ الآتية ليست تحصيلات منطقية في
  $\LogGod[3]$:
  \begin{align*}
    & (p \lif q) \lif (\lnot p \lor q) \\
    & \lnot(\lnot p \land \lnot q) \lif (p \lor q) \\
    & ((p \lif q) \lif p) \lif p \\
    & \lnot(p \lif q) \lif (p \land \lnot q)
  \end{align*}
\end{prob}

\begin{prob}
  أي العلاقات الآتية تصح في منطق غودل؟ أعط جدول صدق لكل منها.
  \begin{enumerate}
    \item $p, p \lif q \Entails q$
    \item $p \lor q, \lnot p \Entails q$
    \item $p \land q \Entails p$
    \item $p \Entails p \land p$
    \item $p \Entails p \lor q$
  \end{enumerate}
\end{prob}

\end{document}
